% Pursuing Stacks (cotes 134-2 to 134-8): archive cote 134-5, site page 112 (PDF page 113), typescript page 299.
% typescript. Draft compared: ps5.tex:373-434.
% AI-transcribed by Claude Opus 5.5 (claude-opus-5-5, adaptive thinking, effort high), Message Batches API, batch ps2 (msgbatch_01NVwr1Yuo8qJf5rykeGBXPK), 2026-11-12; not yet checked by eye.
\scanpage{134-5/112}
\mspage{299}
commutes with sums.  This now is readily expressed by the relations
\[ \varphi'(\underline{I}_A) \simeq \underline{I}_B \quad , \]
\[ \varphi'(\underline{I}_A \mathbin{\mathrm{x}} X) \over{\simeq}{\xrightarrow{\sim}} \xout{\varphi} \varphi'(\underline{I}_A)\mathbin{\mathrm{x}}\varphi'(X) \qquad , \]
i.e.  while $\varphi'$  does \uline{not} commute to finite products in general,  however
it \uline{does} commute to the products which enter in the description of elemen-%
tary homotopies (as these products can be expressed in terms of direct
sums in \xout{M,xM'xxxxxandx\unsure{$\varphi$'}x} A, B ,  and f commutes to these).

Th\over{i}{e}s\hand{e} \inserted{reflexions} suggest\deleted{s} that the notions of homotopy interval structures and
contractibility structures may be generalized,  in a way that the underly-%
ing category need no longer  be stable under finite products nor even ad-%
mit a final object;  and likewise,  the notion of a mo\hand{r}phism of such struc-%
tures may be generalized,  without assuming that the underlying functor
should commute with finite products. The thought that this kind of genera-%
lization may be needed had already occured before\inserted{in these notes}, in connection with
the corresponding situation about twenty five years ago,  when the notion
of a site was developped.  But at present the extension does'nt seem ur-%
gent yet,  and I better stop \xout{this} here this long digression !

The remaining questions  possibly to deal with in part IV  are all
concerned with modelizers.   I'll try to be brief\hand{!}

5)   Consider an ``algebraic structure type'',  and the category \inserted{(M)} of  it's
set-theoretic realizations.  I am looking \over{\unsure{on}}{for} a comprehensive set of suffi-%
cient conditions on M to insure that M is a``canonical modelizer''.   It seems
natural to assume beforehand that in M  (where at any rate small direct and
inverse limits must exist) internal \uline{Hom}'s exist,  and more generally,
for  X,Y two objects over an object S of M,  that $\underline{\mathrm{Hom}}_S(X,Y)$  exists  -
this implies that base change $S'\to S$ in M commutes with small direct limits
and a fortiori,  that direct sums are universal  -  we may as well suppose
them disjoint\inserted{too}.  One feels quite willing\inserted{too} to throw in the total 0-connected-%
ness assumption  (cf section 58),  and that every non-empty object has a
section over the final object.    This preliminary set of conditions on an
algebraic structure is of course highly unusual,  however it is satisfied
for \over{M}{m}ost ``elementary'' algebraic structures\deleted{,} \hand{(}by which I mean $M\hand{\simeq}A\hat{}$ for some
small category A\hand{)},  as well as for  n-stacks or \hand{$\infty$}-Gr-stacks,  for any n
between \over{\illegible{1}}{0} and $\infty$ .  The hope \over{\illegible{1}}{now} is,  in terms of these assumptions,  to
give a necessary and sufficient condition in order that  a)  the ``canonical''
homotopy structure  on M  \over{\unsure{is}}{be} a contractibility structure,  and moreover
b)  the latter structure \over{\unsure{is}}{be} ``modelizing'',  by which we mean that the

\margin{112}
